\chapter{Conclusion and Future Work}
\label{ch:conclusion}

This final chapter summarises the contributions of the project
against the five objectives defined in Chapter~\ref{ch:introduction},
acknowledges the limitations of the current system, and outlines
four concrete directions for future work.

\section{Summary of Contributions}
\label{sec:conclusion-contrib}

The project as presented in this report has delivered five concrete
contributions, each aligned with one of the five objectives from
Section~\ref{sec:objectives}:

\begin{enumerate}
  \item \textbf{A four-class driver stress classifier based on a
  bidirectional LSTM with additive attention.} The classifier
  operates on a five-feature, 30-timestep window sampled at 4~Hz
  and uses only wrist-accessible signals. Its design follows the
  bidirectional-LSTM blueprint of \citet{graves2005bilstm} and the
  attention intuition of \citet{vaswani2017attention}, scaled down
  to a 160{,}000-parameter compact form factor.

  \item \textbf{Empirical accuracy of 93.1\% and macro-F1 of 0.914}
  on the held-out WESAD test split \citep{schmidt2018wesad}, with
  a deployed TensorFlow Lite \citep{tflite2023} footprint of
  approximately 70~KB. This exceeds both the accuracy target
  ($\ge 90\%$) and the size target ($\le 80$~KB) set in
  Chapter~\ref{ch:introduction}.

  \item \textbf{A Wear~OS Kotlin application that performs real-time
  on-device inference.} The application collects wrist sensors at
  their native rates, resamples them to a common 4~Hz grid, applies
  the cached per-feature normalisation learned at training time,
  runs the quantised TensorFlow Lite interpreter, and emits a
  smoothed stress estimate every thirty seconds with end-to-end
  latency well below the 50~ms-per-window target.

  \item \textbf{A documented twelve-byte BLE GATT protocol}
  (service \texttt{ABCD1234\ldots}, characteristic
  \texttt{ABCD1235\ldots}) that carries a timestamp, a 4-valued
  stress level, and a confidence. The protocol is fully specified
  in Section~\ref{sec:m-ble} and can be implemented by any BLE
  central or peripheral independently of the reference
  implementation.

  \item \textbf{A graduated, physiology-driven adaptive cabin
  control system} implemented on the Arduino Nano 33 BLE, comprising
  four active cabin profiles (Calm, Mild, Moderate, Critical), a
  safe \texttt{WAITING} state, and a sixty-second BLE timeout.
  The fail-safe behaviour mirrors the design philosophy of the
  AlcoWatch BAC system under provisional patent ACN1408
  \citep{acn1408}.
\end{enumerate}

Taken together, these five contributions form a complete
end-to-end prototype: a driver wearing a Wear~OS device with this
application installed, while operating a vehicle fitted with the
reference Arduino module, will see their cabin lighting, ventilation,
and audio adapt automatically and proportionally to their
autonomic-nervous-system state, with no cloud dependency and no
compromise of sensor data privacy.

\section{Limitations}
\label{sec:conclusion-limits}

Four limitations must be acknowledged. Addressing them is what the
future-work programme of Section~\ref{sec:conclusion-future} targets.

\textbf{Laboratory versus naturalistic stressors.} The WESAD
dataset \citep{schmidt2018wesad} was collected under controlled
laboratory conditions: subjects were stationary, their posture was
constrained, and stressors were delivered in a fixed protocol (a
variant of the Trier Social Stress Test). Driving stress, by
contrast, is mixed with physical workload (steering, pedal control),
sustained mental workload, and ambient vibration. The physiological
signatures of driving stress may therefore differ in ways that
would reduce the classifier's effective accuracy in deployment.

\textbf{No real-world deployment test.} The system as presented has
not been evaluated \emph{in vivo} on real drivers under real
driving conditions. Our evaluation is limited to held-out WESAD
data and scripted BLE simulations. Before any safety claim can be
made, a supervised on-road study is required.

\textbf{Synthetic stress versus driving stress.} The ``Critical''
class is defined from a per-subject quantile split of the WESAD
stress condition, not from a real driving emergency. It is
therefore an approximation to the autonomic response that a genuine
near-miss event would produce. Whether that approximation is tight
enough for the Critical cabin profile to be triggered at the right
moments in the real world is an empirical question.

\textbf{Accelerometer magnitude discards direction.} We input the
magnitude of the three-axis accelerometer vector as a single
scalar feature. This keeps the feature count to five but discards
directional information: the model cannot distinguish, for example,
between a sharp lateral swerve and a vertical pothole impact, both
of which produce similar magnitudes but very different risk
profiles. Re-admitting the three raw axes as separate features
would grow the input to seven features and slightly increase the
model size, but would likely improve robustness.

\section{Future Work}
\label{sec:conclusion-future}

Four directions suggest themselves for follow-on work.

\textbf{(1) Unification with AlcoWatch into a single platform.} The
AlcoWatch BAC detection system \citep{acn1408} and the present
driver-stress system share the bulk of their architecture: the same
wrist sensors, the same on-device TensorFlow Lite runtime, the same
Arduino-based vehicle controller, and the same fail-safe default-deny
philosophy. A natural next step is to merge the two into a single
Wear~OS application that runs both inference engines concurrently
and exposes two BLE characteristics on one shared service. The
Arduino firmware would then implement a combined decision logic in
which ignition is blocked (AlcoWatch-style) only for BAC above the
legal limit, while the cabin environment adapts continuously
(stress-style) to the driver's autonomic state. This unification
would amortise the hardware cost of the wearable across two distinct
safety outcomes.

\textbf{(2) Real-world data collection during driving.} A follow-on
project should collect a WESAD-style multimodal dataset under
naturalistic driving conditions, either on a driving simulator or,
subject to ethical approval, on a test track. The resulting dataset
would allow the classifier to be fine-tuned on the exact signal
distribution it will encounter in deployment and would provide the
empirical ground for a quantitative safety claim.

\textbf{(3) Personalisation via per-driver calibration.} The
per-subject quantile logic used in the stress-split label mapping
(Section~\ref{sec:m-labels}) is evidence that between-subject
differences in baseline autonomic tone are large enough to matter.
A personalisation layer -- for example, a five-minute baseline
capture the first time the watch is worn, with the resulting per-user
statistics used to normalise sensor streams before they reach the
model -- would likely improve Critical-class recall on drivers whose
resting EDA or heart rate sits at the extremes of the population
distribution.

\textbf{(4) Knowledge distillation to a sub-25~KB student model.}
The current 70~KB model is comfortably deployable on Wear~OS, but
a smaller model would further reduce inference latency and energy
consumption, and would open the door to deployment on ultra-low-power
wearables and hearables. A straightforward approach is knowledge
distillation: train a smaller student network (for example a
single-direction LSTM or a one-dimensional CNN with under 20{,}000
parameters) to match the soft predictions of the current BiLSTM +
attention teacher. If the student can retain, say, 90\% accuracy
in a 25~KB footprint, the inference cost per day drops to the point
where continuous always-on stress monitoring is feasible on a
wristband-class device.

Each of these four directions extends rather than replaces the
contributions of this report, and collectively they define the path
from a validated laboratory prototype to a deployable, personalised,
multi-purpose wearable safety platform.
